% Challenge dossier: VI.01.C1
\subsection*{Challenge VI.01.C1 --- The vanishing ideal of finite affine space}
\label{chal:vi01-c1}

\paragraph{Challenge.}
Let \(k=\mathbb F_q\).  Prove
\[
I(k^n)=
(x_1^q-x_1,\dots,x_n^q-x_n)
\subseteq k[x_1,\dots,x_n].
\]

\ifdefined\FullProblemDossiers
\paragraph{Why this challenge matters.}
It is the multivariable form of the finite-field phenomenon and gives an exact algebraic
presentation of the set of all \(q^n\) rational points.

\paragraph{Guided hints.}
Reduce any polynomial modulo the displayed generators until its degree in every variable is
less than \(q\).  Then prove by induction on \(n\) that such a reduced polynomial vanishing
on all of \(k^n\) must be zero.

\paragraph{Solution.}
The displayed generators vanish on \(k^n\), so their ideal is contained in \(I(k^n)\).
Repeated division by the monic polynomials \(x_i^q-x_i\) writes every \(f\) as
\[
f=\sum_{i=1}^n h_i(x_i^q-x_i)+r,
\]
where \(\deg_{x_i}r<q\) for every \(i\).  If \(f\) vanishes on \(k^n\), then so does
\(r\).

We prove \(r=0\) by induction on \(n\).  For \(n=1\), a polynomial of degree less than
\(q\) with \(q\) roots is zero.  For the induction step write
\[
r=\sum_{j=0}^{q-1}a_j(x_1,\dots,x_{n-1})x_n^j.
\]
Fixing any \((u_1,\dots,u_{n-1})\in k^{n-1}\) gives a degree-\(<q\) polynomial in \(x_n\)
that vanishes at all \(q\) values, so every coefficient value
\(a_j(u_1,\dots,u_{n-1})\) is zero.  Each \(a_j\) also has degree less than \(q\) in each
remaining variable, so the induction hypothesis gives \(a_j=0\).  Hence \(r=0\) and the
claimed ideal equality follows.

\paragraph{Extensions.}
Deduce that every polynomial function \(k^n\to k\) has a unique representative with degree
less than \(q\) in each variable.

\paragraph{Provenance.}
New challenge extending Problem VI.01.P6.
\fi
